\documentclass[11pt]{article}
\usepackage[a4paper,margin=23mm]{geometry}
\usepackage{amsmath,amssymb,amsthm,booktabs,microtype}
\usepackage[hidelinks]{hyperref}
\newcommand{\E}{\mathbb E}\newcommand{\Pp}{\mathbb P}
\newcommand{\CE}{C_{\rm E}}\newcommand{\ph}{\phi_0}
\newtheorem{theorem}{Theorem}\newtheorem{lemma}[theorem]{Lemma}
\title{A five-count comparison with one proof}
\author{}\date{29 October 2026}
\begin{document}\maketitle

The same argument treats the two conditional support strips and a narrower
strip supplied by geometry. Five neighboring counts retain a definite
crossing loss. Four centered replacements estimate that loss at variance
cost, while one fourth moment controls all omitted translates. The
remaining binomial and normal corrections have order $1/n$. Their final
comparison is a pair of concave quartics, so only four endpoint inequalities
are needed for each row of constants.

Write $\ph=(2\pi)^{-1/2}$ and $\CE=(3+\sqrt{10})\ph/6$.
Let $Y=B+U$, where $B$ is Bernoulli of parameter $p$ and
$\E(U\mid B)=0$. Denote the conditional laws by $U_0,U_1$, and put
\[
q=1-p,\quad v=pq,\quad s_j=\E U_j^2,\quad s=qs_0+ps_1,\quad
\lambda=ns,\quad X=\frac{Y-p}{\sqrt{v+s}},\quad\rho=\E|X|^3.
\]
In every statement below the noise laws are those of this one marginal.

\begin{theorem}\label{thm:rows}
Choose any row of Table~\ref{tab:domains}. Suppose $n\ge N$,
$p\in[p_-,p_+]$, $L\le\lambda\le U$, and
\[
-A\le U_0\le a_0,\qquad -a_1\le U_1\le A,\qquad
a_0,a_1\ge0,\qquad a_0+a_1\le d.
\]
For independent copies $X_1,\ldots,X_n$, and every $|H|\le46/25$,
\begin{equation}\label{eq:conclusion}
\sqrt n\{\Pp(X_1+\cdots+X_n\le H)-\Phi(H/\sqrt n)\}
-\CE\rho<-\lambda/100.
\end{equation}
\end{theorem}

\begin{table}[h]
\centering
\begin{tabular}{lrrrrrr}
\toprule
Row&$N$&$[p_-,p_+]$&$A$&$d$&$L$&$U$\\
\midrule
Compact&700&$[3/8,93/200]$&$3/5$&$1/250$&$1/100$&$7/20$\\
Right&200&$[2/5,53/100]$&$3/5$&$9/125$&$3/100$&$9/25$\\
Left&700&$[1/3,2/5]$&$7/10$&$17/250$&$1/100$&$8/25$\\
\bottomrule
\end{tabular}
\caption{Three specializations of the same criterion.}\label{tab:domains}
\end{table}

The compact row is the natural comparison after the narrower geometric
acquisition. The two other rows remain useful without that acquisition.
In particular the right row now starts at $\lambda=.03$, so its smaller-noise
join needs only the existing all-phase $.03$ comparison. Each displayed
theorem is local: acquiring its support and raw threshold, reflection and
an independent finite base remain separate obligations.

\section{A criterion with explicit parameters}

Here is the common criterion used to prove the theorem. Assume
\[
N\ge42,\quad [p_-,p_+]\subset[1/3,2/3],\quad 0<L\le U,\quad
0\le d\le A,\quad d<r,\quad A+d<1,
\]
where
\[
r=\min(p_-,1-p_+),\quad q_*=1-p_+,\quad
v_-=\min_{p_-\le p\le p_+}p(1-p),\quad
v_+=\max_{p_-\le p\le p_+}p(1-p).
\]
The following constants are fixed throughout the proof:
\[
\eta=2/5,\quad \beta=41/100,\quad
\pi_+=22/7,\quad \phi_-=199/500,\quad\phi_+=2/5.
\]
Choose positive $\kappa,\xi,c_-,\zeta,m,V,j_0,j_1$, and
$\bar c\in[0,1)$, satisfying these conditions:
\begin{align}
&\eta\pi_+^2d\le\frac{1-A-d}{A+d},
\qquad 1-\cos(2\pi x)\ge\kappa x^2\quad(|x|\le A+d),\label{eq:kernelguards}\\
&(46/25)^2(v_++U/N)<\xi^2,\quad \xi<1,\quad
N(r-d)>\xi,\quad c_-r\ge1+\xi,\quad c_-<N,\label{eq:phaseguards}\\
&0<\zeta<\min_{z=p_--4/N,\ p_++4/N}z(1-z),\qquad
0<m<2\phi_-\left(1-\frac{8+1/12}{N\zeta}\right),\label{eq:massguards}\\
&V^2<v_-^3,\qquad
\bar c\ge(\xi-q_*)_+,\qquad 4j_0j_1\ge1/q_*,\label{eq:curvguards}\\
&K_0-K_1U/N>0,\quad
K_0=3r-d-\tfrac32(1-2v_-),\quad
K_1=\frac{3(1-2v_-)}{8v_-}.\label{eq:cubicguard}
\end{align}
All divisions are by positive quantities. The points in
\eqref{eq:massguards} lie in $(0,1)$ because $N\ge42$ and
$[p_-,p_+]\subset[1/3,2/3]$.

Put $\alpha=1-c_-/N$, $b=1+3/N$, $z=\kappa\alpha$, and define
\begin{align}
W(\lambda)=\eta m\left\{
\frac{z\lambda-z^2\lambda^2/2+z^3\lambda^3/6-z^4\lambda^4/24}{2}
-\frac{A^2b\lambda+3b^2\lambda^2}{16}
-\frac{4\pi_+^2\lambda}{Nr}\right\}.\label{eq:W}
\end{align}
The common error polynomials are
\[
T=\lambda/q_*,\quad M=A^2T+3T^2,\quad J(\lambda)=j_1\lambda+j_0,
\quad S_c=\frac T4+cJ(\lambda)+\frac{27M}{256}\left(\frac53+\frac{5c}4\right),
\]
\begin{equation}\label{eq:errors}
E_-=\frac{\beta S_{\bar c}+\phi_+(\xi/8+2/27)}V,\qquad
E_+=\frac{\beta S_0+\phi_+(2/27+\xi\lambda/2)}V .
\end{equation}

\begin{lemma}[Four endpoint criterion]\label{lem:criterion}
Under \eqref{eq:kernelguards}--\eqref{eq:cubicguard}, conclusion
\eqref{eq:conclusion} follows if the four rational quantities
\[
P_\pm(L),P_\pm(U)\quad\hbox{are positive},\qquad
P_\pm(\lambda)=W(\lambda)-E_\pm(\lambda)/N-\lambda/100.
\]
\end{lemma}
The proof below supplies the reduction, rather than assuming that endpoint
checks enclose a continuous domain.

\section{One exact crossing identity}

Set $t=np+H\sqrt{v+s}$, $k=\lfloor t\rfloor$, $u=t-k\in[0,1)$ and
$w=t-np$. By \eqref{eq:phaseguards}, $|w|<\xi$ and
$na_0\le t<n(1-a_1)$. On one auxiliary space take independent copies
of the conditional laws and write
\[
T_j=j+\sum_{l\le j}U_{1l}+\sum_{l>j}U_{0l},\quad
L_j(t)=\E(t-T_j)_+,\quad
V_i=\sum_{l<i}U_{1l}+\sum_{l>i}U_{0l}.
\]
The mesh increases because $T_j-T_{j-1}\ge1-a_0-a_1>0$.
Every $T_j$ has exactly the conditional convolution of the actual iid
sample given $j$ successes. Define
\[
h(x)=\E(x-U_0)_+-\E(x-1-U_1)_+-\Pp(1+U_1\le x),\qquad
g_i=\E h(t-i+1-V_i).
\]
Independence of the omitted coordinate gives
\[
g_i=L_{i-1}(t)-L_i(t)-\Pp(T_i\le t).
\]
Here $L_0=t$, $L_n=0$, and $\Pp(T_0\le t)=1$.
Multiplication by the binomial masses $\pi_i$ and finite summation by parts
therefore yields
\begin{equation}\label{eq:cross}
\Pp\Big(\sum X_i\le H\Big)
=(1-u)F_B(k)+uF_B(k+1)
+\sum_{j=1}^{n-1}(\pi_{j+1}-\pi_j)B_j-\sum_{i=1}^n\pi_i g_i,
\end{equation}
where $B_j=L_j(t)-(t-j)_+\ge0$ by Jensen. To verify the interpolation
term, replace every $L_j$ in
$\pi_0+\pi_1t+\sum_{j=1}^{n-1}(\pi_{j+1}-\pi_j)L_j$
by $(t-j)_+$ and telescope once more.
All CDF events in these formulas are closed. The strict survival events
used next are complementary to precisely these closed events.

\section{Five kernels, four centered replacements}

Let $f(x)=\sin^2(\pi x)$ on $[0,1]$ and zero elsewhere.
It is continuously differentiable, its derivative vanishes at the joins,
and that derivative is Lipschitz with constant $2\pi^2$.
We first prove the common pointwise inequality
\begin{equation}\label{eq:minorant} h(x)\ge\eta f(x)\ge0.\end{equation}
For $x<1-a_1$, $h(x)=\E(x-U_0)_+\ge0$, and for $x\ge0$ Jensen
improves this to $h(x)\ge x$. On $[0,1]$ one has $f(x)\le(5/2)x$:
use $\pi^2<10$ up to $1/4$, the concave tangent
$1/2+\pi(x-1/4)$ on $[1/4,2/5]$, and $f\le1$ thereafter.
The two tangent endpoint inequalities follow from $\pi<10/3$.

For $x\ge1-a_1>a_0$, put $D=1+U_1$ and $L=\E(D-x)_+$.
Put--call parity gives $h(x)=\Pp(D>x)-L$. Since $D\le1+A$,
this is nonnegative: it is at least $(x-A)\Pp(D>x)$ until $x=1+A$
and is zero beyond that point. If $1-a_1\le x\le1$, then $L\ge1-x$ and
\[
h(x)\ge\frac{x-A}{1+A-x}L
\ge\frac{1-A-d}{A+d}(1-x)
\ge\eta\pi^2(1-x)^2\ge\eta f(x).
\]
The comparisons at $x=0,1$ are weak, as both sides may vanish.
In particular all atoms of $D$ are included.

Freeze $V=V_{k+1}$ and let $V^{(j)}$ have the law of $V_{k+1+j}$.
There are $n-1$ summands in each. For $j>0$, the latter law replaces
$j$ copies of $U_0$ in $V$ by $U_1$. Put $f_j(x)=f(u-j-x)$.
The two positive changes regroup as
\begin{align*}
&\E f_1(V^{(1)})+\E f_2(V^{(2)})-\E f_1(V)-\E f_2(V)\\
&=\E(f_1+f_2)(V^{(1)})-\E(f_1+f_2)(V)
 +\E f_2(V^{(2)})-\E f_2(V^{(1)}).
\end{align*}
The negative changes have the same form with the laws interchanged.
The translated bumps have disjoint open supports and zero derivatives at
every join. Thus each function in these four replacements still has
derivative Lipschitz constant $2\pi^2$.

For a centered $Z$, Taylor's integral remainder gives
$|\E F(a+Z)-F(a)|\le\frac12\operatorname{Lip}(F')\E Z^2$.
Condition on the other independent summands and apply this twice to each
replacement. The total cost is at most $4\pi^2(s_0+s_1)$, so
\[
\sum_{i=k-1}^{k+3}g_i
\ge\eta\left\{\sum_{j=-2}^2\E f(u-j-V)-4\pi^2(s_0+s_1)\right\}.
\]
The five indices are interior: $|i-np|<\xi+3<4$.
No independence of sums at different counts is used.

The infinite translates of $f$ sum to $\sin^2(\pi(u-V))$.
An omitted nonzero translate with $|j|\ge3$ requires $|V|>2$.
The bump vanishes at both endpoints, so this remains true at collisions.
Consequently
\[
\sum_{j=-2}^2\E f(u-j-V)
\ge \E\sin^2(\pi(u-V))-\E V^4/16.
\]
Writing $\delta=k-np\in(-1-\xi,\xi)$, the actual frozen variance is
\[
\tau=(nq-\delta-1)s_0+(np+\delta)s_1.
\]
The phase guards and $r\ge1/3$ give
\[
(1-c_-/n)\lambda\le\tau\le(1+3/n)\lambda,\quad
s_0+s_1\le\lambda/(nr),\quad
\E V^4\le A^2\tau+3\tau^2.
\]
The last bound follows by expanding the fourth moment of the independent
centered summands, using $\E U_j^4\le A^2s_j$.

For either noise $Z$, an independent copy $Z'$ satisfies
$|Z-Z'|\le A+d$. The cosine guard implies
\[
1-|\E e^{2\pi iZ}|^2
=\E\{1-\cos(2\pi(Z-Z'))\}\ge2\kappa\operatorname{Var}Z.
\]
Hence $|\E e^{2\pi iZ}|\le e^{-\kappa\operatorname{Var}Z}$.
Multiplying only over the independent frozen summands proves
\[
\E\sin^2(\pi(u-V))\ge(1-e^{-\kappa\tau})/2.
\]
We have obtained the generic five-count reserve
\begin{equation}\label{eq:reserve}
\sum_{i=k-1}^{k+3}g_i\ge\eta\left\{
\frac{1-e^{-\kappa(1-c_-/n)\lambda}}2
-\frac{A^2(1+3/n)\lambda+3(1+3/n)^2\lambda^2}{16}
-\frac{4\pi^2\lambda}{nr}\right\}.
\end{equation}

\section{Finite binomial weights and curvature}

We use finite Stirling with its explicit remainder:
\[
\log(j!)=(j+\tfrac12)\log j-j+\tfrac12\log(2\pi)+r_j,\qquad
0<r_j<1/(12j).
\]
For clarity, if $x=1/(2j+1)$ then
$r_j-r_{j+1}=\operatorname{arctanh}(x)/x-1
=\sum_{\ell\ge1}x^{2\ell}/(2\ell+1)$ lies strictly between
zero and $1/(12j(j+1))$. The differences are summable, and Wallis's
product fixes the resulting constant, so $r_j\to0$. Summing the
differences from $j$ to infinity gives the displayed bound. No local central limit theorem
or law-dependent asymptotic remainder is used.

For $k-1\le i\le k+3$ and $z=i/n$, the segment between $z$ and $p$
has variance product at least $\zeta$. Binary relative entropy therefore
satisfies $nD(z\Vert p)\le(i-np)^2/(2n\zeta)<8/(n\zeta)$.
The factorial remainder is bounded below by $-1/(12n\zeta)$.
Since $z(1-z)\le1/4$,
\[
\sqrt n\,\pi_i>
2\phi_-\left(1-\frac{8+1/12}{N\zeta}\right)>m.
\]
All other $g_i$ are nonnegative. The exact reserve \eqref{eq:reserve}
increases with $n$ at fixed $\lambda$. First replace $n$ by $N$ there,
then use
$1-e^{-x}\ge x-x^2/2+x^3/6-x^4/24$ for $x\ge0$.
This proves
\begin{equation}\label{eq:weighted}\sqrt n\sum_i\pi_i g_i\ge W(\lambda).\end{equation}
The order of these steps does not assume monotonicity of a quartic
approximation in $n$.

We give the full summable curvature estimate once. Put
$c=(np-q-t)_+$. If $T,M$ bound $\E W_j^2,\E W_j^4$ on the increasing
side of the binomial mass, where $W_j=T_j-j$, then
\begin{equation}\label{eq:curvature}
\sqrt n\sum_{j=1}^{n-1}(\pi_{j+1}-\pi_j)B_j
\le\frac{\beta}{nv^{3/2}}
\left\{\frac T4+c\sqrt T+\frac{27M}{256}
\left(\frac53+\frac{5c}4\right)\right\}.
\end{equation}
Here is the derivation, including the infinite tail.
For $a=t-j$, centering gives
$B_j=(\E|a-W_j|-|a|)/2\ge0$,
$B_j\le\sqrt T/2$ and $a_+B_j\le T/4$.
The factorization
\[
27x^4-256x+256=(3x-4)^2(3x^2+8x+16)
\]
gives $B_j\le27M/(256|a|^3)$ when $a\ne0$.
The exact adjacent-mass identity is
\[
\pi_{j+1}-\pi_j
=\frac{np-q-j}{(n+1)v}\,
\Pp\{\operatorname{Bin}(n+1,p)=j+1\}.
\]
Discard only its nonpositive terms. The positive numerator is at most
$(a+c)_+$. Since $c\le\bar c<1$, at most one mesh point lies in each of
$(-c,0)$ and $[0,1)$; together they cost $c\sqrt T+T/4$.
For the rest use
$\sum_{i\ge1}i^{-2}<5/3$ and $\sum_{i\ge1}i^{-3}<5/4$.
The first follows from the first five terms and an integral from 5;
the second from the first two terms and an integral from 2.

To finish \eqref{eq:curvature}, the required mode bound holds uniformly
on the entire interval $p\in[1/3,2/3]$:
\[
\max_j\Pp\{\operatorname{Bin}(M,p)=j\}<\frac{\beta}{\sqrt{Mv}}
\quad(M\ge43).
\]
Indeed a mode has $z=j/M=p+\delta/M$, $|\delta|\le2/3$ and
$|1-2p|\le1/3$, while $v\ge2/9$. Thus
$z(1-z)\ge v(1-1/M-2/M^2)>39v/40$.
Finite Stirling and the nonpositive entropy exponent bound the mass by
$(2/5)\sqrt{40/39}(516/515)/\sqrt{Mv}<\beta/\sqrt{Mv}$.
Here $e^{1/516}<516/515$, and the last comparison is a rational square.
Finally $\sqrt n/(n+1)^{3/2}\le1/n$ proves \eqref{eq:curvature}.

A positive adjacent mass requires $j<np$. Its actual variance
$(n-j)s_0+js_1$ is at most $\lambda/q\le T=\lambda/q_*$.
Independence supplies $M=A^2T+3T^2$.
The affine guard gives $J(\lambda)\ge2\sqrt{j_0j_1\lambda}
\ge\sqrt T$, so all these quantities match \eqref{eq:errors}.

\section{Both normal signs and the cubic correction}

Schulz's strict finite-binomial theorem \cite[Theorem 1]{Schulz} states,
for every $0<p<1$, finite $n$ and real argument, that
\[
\sqrt n\{F_B(j)-\Phi((j-np)/\sqrt{nv})\}<\CE\rho_B,
\qquad \rho_B=(1-2v)/\sqrt v.
\]
Apply it at $k,k+1$ and interpolate. No sharper probability-restricted
Bernoulli reserve is used. For $G(x)=\Phi((x-np)/\sqrt{nv})$,
\[
G''(x)\le\frac{\ph(np-x)_+}{(nv)^{3/2}}.
\]
The linear interpolation remainder on the whole unit interval is therefore
\[
\sqrt n\{(1-u)G(k)+uG(k+1)-G(t)\}
\le\frac{\ph}{2nv^{3/2}}u(1-u)(u-w)_+.
\]
For $H\le0$, changing the normal variance from $v$ to $v+s$ is
favorable. Using $u(1-u)\le1/4$ and $u^2(1-u)\le4/27$ gives a total
normal cost at most $\phi_+(\xi/8+2/27)/(nv^{3/2})$.
Here $c=(-q-w)_+\le\bar c$.
For $H\ge0$, $c=0$, interpolation costs at most
$\phi_+(2/27)/(nv^{3/2})$, and the variance correction is at most
$\ph ws/(2v^{3/2})$. The inequality
$1-(1+x)^{-1/2}\le x/2$ proves the latter estimate.
These are exactly the two normal terms in \eqref{eq:errors}.

The conditional supports have fixed signs relative to $p$, since $d<r$.
Expanding their cubes gives the actual raw third absolute moment
\[
v(1-2v)+3v(s_0+s_1)-q\E U_0^3+p\E U_1^3.
\]
The inward bounds $\E U_0^3\le a_0s_0$, $\E U_1^3\ge-a_1s_1$
give an increase of at least $(3r-d)s$. Since
$(1+x)^{3/2}\le1+3x/2+3x^2/8$ for $x\ge0$,
\[
(v+s)^{3/2}(\rho-\rho_B)
\ge s\left\{3r-d-\frac32(1-2v)-\frac{3(1-2v)s}{8v}\right\}
\ge s(K_0-K_1s)>0.
\]
Thus the cubic correction is favorable throughout the stated noise interval.

Combine the exact identity, \eqref{eq:weighted}, \eqref{eq:curvature}
and the two signed normal estimates. The full signed excess is strictly
less than $-W(\lambda)+E_\pm(\lambda)/n$.
Both $E_\pm$ are nonnegative convex quadratics.
Moreover
\[
W''(\lambda)=-\eta m\left\{
\frac{z^2}{4}\bigl[(z\lambda-1)^2+1\bigr]+\frac{3b^2}{8}\right\}<0.
\]
Consequently $P_\pm$ are strictly concave on the whole real line.
Their four positive endpoint values therefore prove $P_\pm>0$ throughout
$[L,U]$. Since $n\ge N$ and $E_\pm\ge0$, this proves
Lemma~\ref{lem:criterion}, including both signs of $H$, every phase and
all threshold atoms.

\section{Checking the rows}

The remaining parameters are in Table~\ref{tab:parameters}; all omitted
quantities are defined by the domain row, not chosen independently.
\begin{table}[h]
\centering\small
\begin{tabular}{lrrrrrrrrr}
\toprule
Row&$\kappa$&$\xi$&$c_-$&$\zeta$&$m$&$V$&$\bar c$&$j_0$&$j_1$\\
\midrule
Compact&3&$93/100$&6&$23/100$&$3/4$&$113/1000$&$79/200$&$1/5$&$5/2$\\
Right&3&$231/250$&5&$23/100$&$13/20$&$117/1000$&$23/50$&$9/40$&$5/2$\\
Left&$7/5$&$23/25$&6&$11/50$&$3/4$&$13/125$&$8/25$&$13/200$&$250/39$\\
\bottomrule
\end{tabular}
\caption{Constants for the common analytic criterion.}\label{tab:parameters}
\end{table}

Here is an elementary check of the only trigonometric guard.
For Compact and Right, use $\theta=2/3$ and
$\sin(\pi\theta)=\sqrt3/2>17/20$. On $[0,\theta]$, concavity gives
$\sin(\pi x)\ge(51/40)x$, whose squared coefficient proves $\kappa=3$.
The Compact diameter $A+d$ is below $\theta$.
For Right, on $[\theta,84/125]$ the sine is at least
$17/20-(22/7)(2/375)>0$; twice its square exceeds
$3(84/125)^2$.
For Left take $\theta=3/4$ and
$\sin(\pi\theta)=\sqrt2/2>7/10$.
The chord slope is at least $14/15$. On
$[\theta,96/125]$ the lower sine bound is $563/875$.
Twice these squared bounds prove $\kappa=7/5$.
At zero the comparisons are weak; evenness handles negative arguments.

Every other guard and endpoint is rational, apart from the fixed normal
constants enclosed by the elementary bounds $25/8<\pi<22/7$.
The checker reconstructs the polynomials directly from
\eqref{eq:W}--\eqref{eq:errors}, verifies the completed-square identity for
$W''$, and checks all guards before the endpoints. Conservative lower
bounds for the endpoint margins are:
\[
\begin{array}{c|rrrr}
&P_-(L)&P_-(U)&P_+(L)&P_+(U)\\ \hline
\text{Compact}&.002314&.061304&.003341&.063680\\
\text{Right}&.000461&.000885&.004654&.011247\\
\text{Left}&.000055&.018585&.000897&.022487
\end{array}
\]
The displayed decimals are exact rational lower targets, not floating-point
proofs. No grid, fitted coefficient array or truncation of an incoming-count
series occurs. This completes the proof of Theorem~\ref{thm:rows}.

\section*{Scope}
This note replaces the repeated crossing-loss, mass, curvature, normal and
cubic arguments in the three v15 comparison notes with one criterion.
It proves its local statements for every actual marginal satisfying a row.
It does not acquire that row for arbitrary laws, turn a signed conclusion
into an absolute one without reflected-law hypotheses, or establish a finite
base. The smaller-noise comparisons and the geometric acquisition are
separate joins. All previous sources and certificates are preserved.

\begin{thebibliography}{1}
\bibitem{Schulz}J. Schulz,
\emph{The optimal Berry--Esseen constant in the binomial case},
doctoral dissertation, Universit\"at Trier, 2016, Theorem 1, printed page 1.
\url{https://d-nb.info/1197702695/34}.
\end{thebibliography}
\end{document}
